%------------------------
% Resume Template
% Author : Anubhav Singh
% Github : https://github.com/xprilion
% License : MIT
%------------------------

\documentclass[a4paper,20pt]{article}

\usepackage{latexsym}
\usepackage[empty]{fullpage}
\usepackage{titlesec}
\usepackage{marvosym}
\usepackage[usenames,dvipsnames]{color}
\usepackage{verbatim}
\usepackage{enumitem}
\usepackage{hyperref}
\usepackage{fancyhdr}
\usepackage[UTF8, fontset=fandol]{ctex}

\pagestyle{fancy}
\fancyhf{} % clear all header and footer fields
\fancyfoot{}
\renewcommand{\headrulewidth}{0pt}
\renewcommand{\footrulewidth}{0pt}

% Adjust margins
\addtolength{\oddsidemargin}{-0.530in}
\addtolength{\evensidemargin}{-0.375in}
\addtolength{\textwidth}{1in}
\addtolength{\topmargin}{-.45in}
\addtolength{\textheight}{1in}

\urlstyle{rm}

\raggedbottom
\raggedright
\setlength{\tabcolsep}{0in}

% Sections formatting
\titleformat{\section}{
  \vspace{-10pt}\scshape\raggedright\large
}{}{0em}{}[\color{black}\titlerule \vspace{-6pt}]

%-------------------------
% Custom commands
\newcommand{\resumeItem}[2]{
  \item\small{
    \textbf{#1}{: #2 \vspace{-2pt}}
  }
}

\newcommand{\resumeItemWithoutTitle}[1]{
  \item\small{
    {\vspace{-2pt}}
  }
}

\newcommand{\resumeSubheading}[4]{
  \vspace{-1pt}\item
    \begin{tabular*}{0.97\textwidth}{l@{\extracolsep{\fill}}r}
      \textbf{#1} & #2 \\
      \textit{#3} & \textit{#4} \\
    \end{tabular*}\vspace{-5pt}
}


\newcommand{\resumeSubItem}[2]{\resumeItem{#1}{#2}\vspace{-3pt}}

\newcommand{\resumeProjectItem}[3]{
  \item
    \begin{tabular*}{0.97\textwidth}{l@{\extracolsep{\fill}}r}
      \textbf{#1} & \textit{#2} \\
    \end{tabular*}
    \vspace{-2pt}
      #3
    \vspace{-2pt}
}

\renewcommand{\labelitemii}{$\circ$}

\newcommand{\resumeSubHeadingListStart}{\begin{itemize}[leftmargin=*]}
\newcommand{\resumeSubHeadingListEnd}{\end{itemize}}
\newcommand{\resumeItemListStart}{\begin{itemize}}
\newcommand{\resumeItemListEnd}{\end{itemize}\vspace{-5pt}}

%-----------------------------
%%%%%%  CV STARTS HERE  %%%%%%
 
\begin{document}

%----------HEADING-----------------
\begin{tabular*}{\textwidth}{l@{\extracolsep{\fill}}r}
   \textbf{{\LARGE 聂家明}} \\
   \href{https://www.linkedin.com/in/jmnie/}{LinkedIn: in/jmnie}  & 邮箱:\href{mailto:jiaming.nie13@gmail.com}{jiaming.nie13@gmail.com} \\
  & 电话: (+86)-135-2171-3097
\end{tabular*}

%-----------EDUCATION-----------------
\section{教育背景}
  \resumeSubHeadingListStart
    \resumeSubheading
      {伍斯特理工学院}{马萨诸塞，美国}
      {硕士 - 计算机科学}{2017.8 - 2019.12}
    %   {\scriptsize \textit{ \footnotesize{\newline{}\textbf{Courses:} Operating Systems, Data Structures, Analysis Of Algorithms, Artificial Intelligence, Machine Learning, Networking, Databases}}}
    \resumeSubheading
      {利物浦大学}{利物浦，英国}
      {学士 - 电气工程}{2015.8 - 2017.6}
      \resumeSubheading
      {西交利物浦大学}{苏州，中国}
      {学士 - 电气工程}{2013.8 - 2015.8}
    \resumeSubHeadingListEnd
	    
\vspace{2pt}
\section{技能}
	\resumeSubHeadingListStart
	\resumeSubItem{语言}{~~~Python, JavaScript/TypeScript, Java}
    \resumeSubItem{AI/LLM}{~~~LLM Fine-Tuning (SFT/RFT), RAG}
    \resumeSubItem{后端/部署}{~~~FastAPI, Flask, Node.js, Redis, MongoDB, Docker, AWS, Prometheus, Grafana}
	\resumeSubHeadingListEnd
\vspace{4pt}
\section{工作经历}
  \resumeSubHeadingListStart
  \vspace{4pt}
  \resumeSubheading{Youlify}{远程}
    {创始软件工程师}{2025.1 - 至今}
    \resumeItemListStart
        \item{独立设计并实现医疗计费代码抽取的端到端 \textbf{LLM} 微调与评测流水线，基于 2,000+ 份临床记录构建 \textbf{SFT/RFT} 数据集，自动化完成数据处理、训练任务编排、检查点跟踪、效果评测与成本分析。}
        \item{建设拒赔审查与申诉信生成的后端工作流，将理赔单、\textbf{EOB}、保险数据等转化为可追溯上下文，完成模型路由、异步任务触发、状态回传和业务状态管理，支撑长链路理赔自动化任务。}
        \item{主导单体后端向 \textbf{FastAPI} 微服务架构迁移，拆分诊所管理、Portal API、鉴权和 \textbf{GenAI} 服务，并搭建 \textbf{Prometheus + Grafana} 可观测体系，提升系统稳定性与迭代效率。}
        \item{设计实时任务基础设施，使用 \textbf{Redis} 结合 \textbf{WebSocket} 为长耗时异步任务推送进度，并利用 \textbf{SSE} 流式传输前端 \textbf{AI} 聊天交互。}
        \item{引入 \textbf{Apache Airflow} 编排患者 \textbf{EHR} 数据与临床文档同步长任务，打通 Portal、Integration 与 Automation Jobs 的任务触发、取消、状态持久化与恢复链路，支撑多诊所、多 \textbf{EHR} 系统的数据同步与进度追踪。}
    \resumeItemListEnd
    
    \vspace{4pt}
  \resumeSubheading{北京峥研软件有限责任公司（peerup）}{上海, 中国}
    {软件开发工程师}{2023.6 - 2025.1}
    \resumeItemListStart
        \item{使用 \textbf{FastAPI} 为笔记产品中的 \textbf{AI} 助手和聊天机器人构建后端服务，支持流式响应、文件感知交互以及在编辑器内由助手触发的操作。}
        \item{开发了基于 \textbf{Milvus} 和 \textbf{Elasticsearch} 的 \textbf{RAG} 流水线，结合公共知识库、用户内容及外部数据源为助手回答提供依据。}
        \item{在 \textbf{FastAPI} 生态中使用 \textbf{LangChain} 与 \textbf{LlamaIndex} 封装检索、上下文组装和外部服务调用能力，支撑多步骤 \textbf{AI} 功能在后端稳定执行。}
      \resumeItemListEnd

    \vspace{4pt}
    \resumeSubheading{Oqton}{上海, 中国}
    {软件开发工程师}{2021.10 - 2023.2}
    \resumeItemListStart
         \item{开发 \textbf{Python} 服务与工业设备进行交互，采集机器数据并通过 \textbf{API} 接口发送控制指令。}
        \item{在 \textbf{Linux} 网关上部署并维护 \textbf{Docker} 应用，协助排查跨 \textbf{Elasticsearch} 和 \textbf{Kafka} 系统的生产环境问题。}
        \item{使用 \textbf{Node.js}、\textbf{Express} 和 \textbf{React.js} 开发后端与前端功能，支持机器配置及数据管理工作流。}
      \resumeItemListEnd

    \vspace{4pt}
    \resumeSubheading
		{微福思软件科技有限公司(Micro Focus)}{上海, 中国}
		{软件开发工程师}{2020.7 - 2021.9}
		\resumeItemListStart
              \item{使用 \textbf{JavaScript}、\textbf{Node.js} 和 \textbf{C++} 开发用于 Web 组件检测的自动化测试功能。}
                \item{实现了跨域组件检测功能，并在产品开发期间优化了内部案例处理工作流。}
        % \resumeItem{Project Course - Find Movie Similarity from Plot Summaries}
        %   {Created project based course using Unsupervised learning and natural language processing.}
        % \resumeItem{Tutorial - Introduction to Reinforcement Learning}
        %   {Created tutorial for Q-learning RL algorithm and  concepts.}
        % \resumeItem{Impact}{Course has been taken by 250+ students so far with 4.65 average rating.}
		\resumeItemListEnd

\resumeSubHeadingListEnd

%---------------PUBLICATIONS-------------
% \vspace{-5pt}
% \section{Publications}
% \resumeSubHeadingListStart
% % \resumeSubItem{Book: Deep Learning on Web (Web Development, Deep Learning)}{Work in Progress book to be published by Packt Publishing in late 2019. Tech: Django, Python, AWS, GCP, Azure (November '18)}
% \item{R Li, GP Balakrishnan, \textbf{J Nie}, Y Li, E Agu, M Stein, A Abrantes, D Herman, On smartphone sensability of bi-phasic user intoxication levels from diverse walk types in standardized field sobriety tests, \textit{IEEE EMBC 2019}, 3279-3285}
% \vspace{2pt}
% \item{R Li, GP Balakrishnan, \textbf{J Nie}, Y Li, E Agu, K Grimone, D Herman, On smartphone sensability of bi-phasic user intoxication levels from diverse walk types in standardized field sobriety tests. \textit{IEEE Access 9}, 61237-61255}
% \resumeSubItem{Book: Deep Learning on Mobile Devices (Flutter App Development, Deep Learning)}{Work in Progress book to be published by Packt Publishing in late 2019. Tech: Flutter, Android, Firebase, TensorFlow, Python, Dart (December '18)}

%-----------PROJECTS-----------------
\vspace{-5pt}
\section{相关项目}
\resumeSubHeadingListStart


\vspace{4pt}
\resumeProjectItem{用于医疗计费代码提取的 LLM 微调工作流}{2026.2 - 至今}{
\resumeItemListStart
    \item{设计并实现了端到端 \textbf{OpenAI} 托管的微调流水线，以实现医疗计费代码提取的自动化。处理了超过 2,000 份临床记录，并将其整理为 \textbf{SFT} 和 \textbf{RFT} 数据集，用于模型的训练、验证和测试。}
    \item{通过 \textbf{OpenAI Fine-Tuning API} 编排 \textbf{GPT-4.1 SFT} 和 \textbf{o4-mini RFT} 训练任务，实现了文件上传管理、状态轮询、检查点追踪和实验元数据采集的自动化。}
    \item{为 \textbf{RFT} 设计了自定义 \textbf{Python} 评分器以优化结构化 \textbf{JSON} 输出，结合加权编辑距离与领域特定惩罚项（\textbf{CPT} 前缀、\textbf{ICD-10} 和修饰符）对预测结果进行评分。}
    \item{建立了全面的评估框架，追踪所有计费代码的宏观/微观指标，并结合 \textbf{Token} 成本估算和检查点分析，与基准推理模型进行对标。}
\resumeItemListEnd
}

\vspace{4pt}
\resumeProjectItem{医疗 SaaS 数据同步与实时任务平台}{2025.1 - 至今}{
\resumeItemListStart
    \item{设计并实现基于 \textbf{Apache Airflow} 的患者 \textbf{EHR} 数据同步编排，将 Athena、Practice Fusion 等传统医疗系统的长耗时拉取任务迁移为可触发、可取消、可追踪的 \textbf{DAG}。}
    \item{使用 \textbf{FastAPI + Gunicorn/Uvicorn} 构建高并发后端 API，并通过 \textbf{WebSocket} 向前端推送抓取进度、错误状态和任务结果。}
    \item{基于 \textbf{Redis} 分布式锁、任务队列状态和 \textbf{MongoDB} 多租户存储设计幂等处理流程，保障诊所租户间的数据隔离与失败重试。}
    \item{配合客户接入和上线排障，沉淀可观测日志、任务状态查询和异常恢复流程，缩短新诊所接入与数据验证周期。}
\resumeItemListEnd
}

\vspace{4pt}
\resumeProjectItem{笔记应用多功能 AI 助手与对话机器人}{2023.9 - 2025.1}{
\resumeItemListStart
    \item 主导开发了笔记应用的 AI 助手后端（基于 \textbf{FastAPI}），支持用户通过快捷键唤起助手。实现了基于 \textbf{SSE} (Server-Sent Events) 的流式响应机制，支持复杂指令执行及文件上传问答。
    \item 构建了基于 \textbf{Milvus} 向量数据库和 \textbf{ElasticSearch} 的混合检索 RAG 系统，整合公共知识库、用户个人知识库及第三方网络搜索结果，显著增强了 LLM 回复的准确性与上下文相关性。
    \item 使用 \textbf{LangChain} 与 \textbf{LlamaIndex} 封装检索、上下文构造和外部服务调用能力，支撑后端多步骤 \textbf{AI} 任务执行。
    \item 设计了基于 \textbf{Redis} 的会话管理机制，高效处理消息、对话及附件 ID 的映射与缓存，并实现了基于用户层级的数据隔离与检索权限控制。
    \item 集成 \textbf{阿里云 OSS} 开发文件上传接口，打通了大语言模型对用户上传文档的深度理解与问答通道。
\resumeItemListEnd
}

\vspace{4pt}
\resumeProjectItem{文档内容增强后端系统}{2023.10 - 2024.1}{
\resumeItemListStart
    \item 开发了基于 \textbf{FastAPI} 的内容增强后端服务，通过智能分词与关键词高亮技术，大幅提升了用户在笔记应用中的阅读与交互体验。
    \item 引入 \textbf{Aho-Corasick} 算法实现高性能文本分词，确保前端能从海量输入中实时、精准地识别关键术语。
    \item 完成了从 MySQL 到 \textbf{ElasticSearch} 的数据迁移与同步，构建了高性能的搜索接口，支持用户点击关键词即时获取详细定义与百科信息。
    \item 设计并维护了 MySQL 中的关键词映射表，利用网络爬虫预处理的数据建立索引，实现了毫秒级的内容关联与详情检索。
\resumeItemListEnd
}

% \vspace{4pt}
% \resumeProjectItem{安全 EHR 数据同步流水线}{2025.1 - 至今}{
% \resumeItemListStart
%     \item 开发了高适应性的数据提取流水线，通过在安全无头浏览器环境中模拟经过身份验证的用户交互，自动同步来自不同第三方 EHR 系统的患者记录 [cite: 46]。
%     \item 使用 \textbf{FastAPI} 构建 API 服务，通过 Gunicorn 配合多个 Uvicorn worker 部署以处理并发请求；实现了 \textbf{WebSocket} 接口以向前端推送实时抓取进度和状态更新 [cite: 47, 48]。
%     \item 设计了基于 \textbf{Redis} 的多租户状态管理系统以缓存同步进度，并使用 \textbf{MongoDB} 进行统一数据存储，确保了租户数据流之间的隔离 [cite: 49]。
%     \item 利用 \textbf{Playwright} 实现健壮的浏览器自动化，使用 \textbf{Redis} 进行任务队列管理和分布式锁控制，确保了可靠且幂等的数据处理，显著缩短了客户上线前的技术准备时间 [cite: 50]。
% \resumeItemListEnd
% }

% \vspace{4pt}
% \resumeProjectItem{产品笔记页的多功能AI助手和对话机器人}{2023.10 - 至今}{
% % \vspace{-2pt}
%  \resumeItemListStart
%  \item 开发了基于 \textbf{FastAPI} 笔记页 AI 助手和对话机器人后端，用户可在笔记页（类似 OneNote 和 Notion）按空格调出助手，通过提问获取流式（Server-Send Event) 输出的回复并执行操作，同时上传文件并对内容进行提问。使用\textbf{MySQL}和\textbf{Elastic Search}进行数据存储，\textbf{redis} 作为缓存中间件。
% \item 构建了基于\textbf{Milvus}向量数据库和\textbf{ElasticSearch}的 RAG 系统，用于召回公共知识库，用户个人知识库以及第三方网络检索信息，加载到请求大语言模型API的上下文中，增强对用户的回答。
% \item 基于\textbf{LangChain}和\textbf{llamaindex} 封装后端检索、上下文构造和外部服务调用能力，用于其他服务调用 API 时执行相关 task。
% \item 使用\textbf{Redis}对消息、对话和附件 ID 进行管理，同时利用 Redis 缓存用户信息（通过访问后端API获取），并基于用户层级进行相应的数据检索。
% \item 项目中中基于阿里云OSS开发了文件上传接口，并支持用户针对文件进行提问。
% \resumeItemListEnd
% }

% \vspace{4pt}
% \resumeProjectItem{笔记页产品的内容增强项目（知识库）}{2023.10 - 2024.2}{
% % \vspace{-2pt}
% \resumeItemListStart
%  \item 开发了一个基于\textbf{FastAPI}的后端，提升记笔记应用中的用户体验。当用户输入文本时，后端执行分词，使前端能够高亮显示关键词。用户可以点击这些关键词以检索每个术语的详细信息，丰富应用内的内容互动。
%  \item 使用\textbf{Aho-Corasick}算法实现文本分词，以高效地从前端输入中识别关键词。
%  \item 设计了\textbf{RESTful API}来执行文本拆分和检索用户从前端输入的关键词信息。
%  \item 将\textbf{ElasticSearch}集成到后端，将内容从\textbf{MySQL}数据库转移到ElasticSearch以支持搜索接口获取数据。
% \item 在MySQL表格中创建关键词ID映射表，将爬虫爬取的数据做成对应关系，以便前端快速获取相应内容。
% \resumeItemListEnd
% }

% \vspace{4pt}
% \resumeProjectItem{集成硬件机器的全栈项目}{2021.10-2022.10}{
% % \vspace{-2pt}
% \resumeItemListStart
% \item 在Oqton期间一系列的客户项目，用于开发基于\textbf{Python}和\textbf{MQTT}的IoT项目，将客户机器接入产品后端。
% \item 开发Python应用程序，使用Flask框架与机器API进行交互，收集机器数据并向机器发送命令。
% \item 使用Docker在Linux网关上部署Python应用程序，并进行维护以确保数据传输正确。
% \item 解决由DevOps团队提出的故障单，并使用Elastic Search和Apache Kafka定位问题。
% \item 使用Node.js开发RESTful API并创建新功能，以对不同机器的替代字段进行过滤和排序。
% \item 使用React.js开发前端功能，连接这些机器并配置替代页面。
% \item 实施和维护Python存储库的CI/CD流水线，从而提高单元测试覆盖率和整体软件可靠性。
% \resumeItemListEnd
% }

% \vspace{4pt}
% \resumeProjectItem{Web产品用户界面端自动化测试项目}{2022.10 - 2023.2}{
%     % \vspace{-2pt}
%      \resumeItemListStart
%      \item 使用Selenium框架(JavaScript)开发和维护自动化测试框架，执行产品的自动化测试计划。
%      \item 借助已有的测试用例，开发并拓展相应更复杂的场景并持续优化已有的代码逻辑和继承关系。
%      \item 与前端开发人员合作，将API测试集成到UI测试框架中，实现对Canvas渲染功能的完整测试。
%      \item 使用Docker运行Selenium脚本，基于CircleCI完善开发相应的CI/CD流水线，周期性运行测试脚本并收集数据。
%      \item 根据测试收集的结果反馈给开发团队，根据测试结果提高自动化测试覆盖率。
%     \resumeItemListEnd
% }

% \vspace{4pt}
% \resumeProjectItem{基于JavaScript和Django的微信小程序项目}{2021.9 - 2021.12}{
%     % \vspace{-2pt}
%      \resumeItemListStart
%      \item 在业余时间和视觉设计师合作开发的微信小程序项目，用于展览等不同场景。
%     \item 基于微信小程序开发平台，使用原生JavaScript和相应的微信API，python Django作为后端，构建了一个网站和一个微信小程序平台。
%     \item 使用阿里云部署服务后端，为前端提供暴露相应的API，并使用MySQL存储相应数据。
%     \item 用于展览的场景是参展者会在一个使用投影仪放映二维码的空间里，使用微信小程序扫描并解锁相应的徽章。
%     \item 开发了基于对话交互式的微信小程序，用户可以通过选择并得到不同的故事线，最终解锁相应剧情结果。
%     \resumeItemListEnd
% }

% \vspace{6pt}
% \resumeSubItem{人体运动数据与醉酒程度的研究}{
%     \resumeItemListStart
%     \item{为期两年的实验室项目，根据移动设备收集人体运动时的加速度计和陀螺仪数据，利用深度学习对数据进行分
% 类，以表明人体的醉酒程度。}
%     \item 以共同作者身份发表了题为“On Smartphone Sensibility of Bi-Phasic User Intoxication Levels from Diverse Walk Types in Standardized Field Sobriety Tests”，并在2019年的IEEE EMBC会议上进行了展示。
%     \item 通过引入多线程编程，优化了Python数据处理脚本，处理时间减少了70\%。
%     \item 研究了多种深度学习网络（CNN和RNN）以预测步行数据中的醉酒水平，并将预测准确率从50％提高到80％。
%     \item 将LSTM回归模型部署在Android应用程序上，用于实时监测用户的醉酒水平。
%       \resumeItemListEnd}


% \resumeSubItem{Inferring Intoxication Level From Human Motion Data 
% }
% \vspace{2pt}
% \resumeSubItem{Reinforcement Learning based Traffic Control System (Reinforcement Learning, Computer Vision)}{AI model to resolve city traffic around 50\%
% faster. Tech: Python, Alibaba Cloud, Raspberry Pi, Arduino, SUMO \& OpenCV. (August '18)}

% \vspace{2pt}
% \resumeSubItem{使用深度学习的风格转换 (Style Transfer) }{
% \begin{itemize}
% \item 开发了一个实时深度学习算法，用于重构具有多个输入风格的图像。
% \item 设计和实现了一种新颖的体系结构，使用两个独立的深度神经网络从风格图像中提取特征并生成目标图像。
% \item 实现了任何尺寸图像的实时多风格转移，显著提高了以前的性能基准。
% \item 使用Apache MXNet实现了该项目，并创建了一个可扩展的框架，用于将新的风格图像添加到模型中。   
% \end{itemize}
% \resumeItemListEnd}

% \vspace{2pt}
% \resumeSubItem{模拟操作系统项目 (C语言)}{
% \begin{itemize}
% \item 该项目使用纯C语言开发。
% \item 具备全面的进程管理机制，可创建和销毁进程。
% \item 多线程机制使多个任务能够同时运行，提高效率。
% \item 通过先进的磁盘IO调度系统和其他相关功能管理磁盘读写操作。
% \item 文件系统采用索引和缓存设计以提高性能并降低延迟。
% \item 内存管理系统采用交换机制和虚拟分页以优化内存使用。
% \item 另外，还实现了多定时器管理，提供精确的时间控制。
% \end{itemize}
% }

% \vspace{2pt}
% \resumeSubItem{集成硬件机器的全栈项目}{
% \begin{itemize}
% \item 在Oqton（2021.10-2023.2）期间，我全面负责了一系列项目，包括：
% \item 开发Python应用程序，使用Flask框架与机器API进行交互，收集机器数据并向机器发送命令。
% \item 使用Docker在Linux网关上部署Python应用程序，并进行维护以确保数据传输正确。
% \item 解决由DevOps团队提出的故障单，并使用Elastic Search和Apache Kafka定位问题。
% \item 使用Node.js开发RESTful API并创建新功能，以对不同机器的替代字段进行过滤和排序。
% \item 使用React.js开发前端功能，连接这些机器并配置替代页面。
% \item 实施和维护Python存储库的CI/CD流水线，从而提高单元测试覆盖率和整体软件可靠性。
% \end{itemize}
% }

% \resumeSubItem{Full-Stack Web Application Games (Web Development, Vue)}{Created Delivered several full-stack projects collaborating with visual designers in spare time. It contains several web application in the format of website and WeChat Mini-program Platform for different purposes.}

% \resumeSubHeadingListEnd
% \begin{itemize}
% \item{Images clicked using drones, provided by ISRO were stitched together using distributed public compute nodes, effectively bringing down processing time exponentially. Tech: PHP, C++, Java, Python (March '18)}
% \end{itemize}
% \vspace{2pt}
% \resumeSubItem{Drag-n-drop machine learning learning environment (Web Development, Machine Learning)}{Scratch like tool for implementing machine learning pipelines along with built in tutorial for each concept. Tech: Python, JavaScript (September '18)}
% \vspace{2pt}
% \resumeSubItem{Search Engine and Social Network(Web Development, Web Crawler, Search)}{Created from scratch a social network and a search engine based on the idea of integrating Facebook and Google. The launched website was among top 1000 websites in India during 2012-2013. Tech: PHP, MySQL, HTML, CSS, WebSockets, JavaScript, RSS, XML ( May '12)}

%\vspace{-5pt}
%-----------Awards-----------------
% \section{Honors and Awards}
% \begin{description}[font=$\bullet$]
% \item {Awarded title of Intel Software Innovator - May, 2019}
% \vspace{-5pt}
% \item {Second Runner's Up at TCS EngiNx Engineering Project Innovation Content - September, 2018 }
% \vspace{-5pt}
% \item {Runner's Up at Facebook Developers Circle Hackathon - August, 2017}

%\end{description}

% \vspace{-5pt}
% \section{Volunteer Experience}
%   \resumeSubHeadingListStart
% 	\resumeSubheading
%     {Community Lead at Developer Student Clubs NSEC}{Kolkata, India}
%     {Conducted online and offline technical \& soft-skills training impacting over 3000 students.}{Jan 2019 - Present}
% \vspace{5pt}
%     % \vspace{10pt}\textbf{\large{Community Experience}}
%     \resumeSubheading
%     {Event Organizer at Google Developers Group Kolkata}{Kolkata, India}
%     {Organized events, conducted workshops and delivered workshops reaching over 7000 developers.}{Jan 2018 - Present}

\resumeSubHeadingListEnd

\end{document}
